% TEMPLATE-MILC-v3.tex - plantilla maestra UNICA de las guias MILC v3.
%
% La usan las tres familias de guias, cada una con su driver:
%   · Grados 6-11  -> scripts/build-guias-g11.py
%   · Bebras       -> scripts/build-guias-bebras.py
%   · Semillero    -> scripts/build-guias-semillero.py
%
% Antes habia tres copias casi identicas (~400 lineas iguales, 6 distintas):
% cada mejora habia que aplicarla tres veces, y en la practica no se hacia
% (el motor de recursos vivia solo en dos de ellas). Lo que varia entre
% familias esta parametrizado en 6 placeholders que cada driver llena
% (se nombran aqui SIN su sintaxis real: el builder reemplaza por texto plano,
% tambien dentro de los comentarios, y un valor multilinea rompe el preambulo):
%
%   HEADER_IZQ        encabezado izquierdo de pagina
%   HEADER_DER        encabezado derecho de pagina
%   PORTADA_ETIQUETA  rotulo sobre el numero grande (GRADO/MOMENTO/MODULO)
%   PORTADA_SERIE     linea de serie de la portada
%   PORTADA_ID        identificador de la guia en la portada
%   PORTADA_CIRCULO   el circulito de grado (°), o vacio si no aplica
%
% El resto de claves las documenta content/guias/_SCHEMA.md. El script valida
% que no queden placeholders sin reemplazar antes de escribir el archivo final.

\documentclass[11pt,letterpaper]{article}
\usepackage[margin=2.0cm]{geometry}
\usepackage[table]{xcolor}
\usepackage{fontspec}
\usepackage[spanish]{babel}
\usepackage{tikz}
% Librerías TikZ para los diagramas de las guías (bloque `recursos:`):
%  · babel  → babel en español vuelve activos caracteres como «>», y TikZ los usa
%             en sus opciones (>=stealth, ->). Sin esta librería, cualquier
%             diagrama con flechas revienta con «\language@active@arg> has an
%             extra }». Debe ir ANTES de que se dibuje nada.
%  · positioning        → sintaxis «below left=2pt and 3pt».
%  · decorations.pathreplacing → llaves ({brace}) para señalar rangos.
\usetikzlibrary{babel,positioning,decorations.pathreplacing}
\usepackage{graphicx}
% Circuitos: el trabajo de electrónica se hace hoy con micro:bit y sus sensores
% integrados (MakeCode, sin cableado), así que no se necesitan esquemáticos.
% Si algún día entran montajes con protoboard, basta descomentar esta línea:
% los diagramas .tex/.tikz declarados en `recursos:` ya se incrustan solos.
% \usepackage{circuitikz}
\usepackage[most]{tcolorbox}
\usepackage{tabularx}
\usepackage{array}
\usepackage{fancyhdr}
\usepackage{titlesec}
\usepackage[document]{ragged2e}  % bandera a la derecha en todo el cuerpo
\usepackage{enumitem}
\usepackage{microtype}

% Helvetica Neue no trae la flecha U+2192, asi que cada «→» escrito literalmente
% en el YAML se compilaba a NADA, en silencio (462 flechas en 61 guias: rutas del
% tipo «paleta → tipografia → lockup» salian rotas). Mapeamos el caracter a su
% version matematica, que si tiene glifo. Asi el autor puede seguir escribiendo
% «→» comodamente en el YAML.
\usepackage{newunicodechar}
\newunicodechar{→}{\ensuremath{\rightarrow}}
% Mismo problema con los simbolos de marca y vinetas: el «✓» de «una palomita ✓»
% salia como cuadrito vacio en el PDF de todas las guias que lo usaban.
\usepackage{pifont}
\newunicodechar{✓}{\ding{51}}
\newunicodechar{✗}{\ding{55}}
\newunicodechar{✔}{\ding{52}}
\newunicodechar{•}{\textbullet}
\newunicodechar{≥}{\ensuremath{\geq}}
\newunicodechar{≤}{\ensuremath{\leq}}

\setmainfont{Helvetica Neue}
\setsansfont{Helvetica Neue}
\setlength{\parindent}{0pt}
\setlength{\parskip}{6pt}
% Sin particion de palabras: en 6.º-7.º una palabra cortada por guion al final
% de linea («Comuni-cacion») frena la lectura mucho mas que una linea corta.
\hyphenpenalty=10000
\exhyphenpenalty=10000
\pretolerance=2000
\tolerance=3000
\setlength{\emergencystretch}{3em}
\linespread{1.10}
\renewcommand{\arraystretch}{1.25}
\setlist{leftmargin=5mm,itemsep=1mm,topsep=1mm}
\newcolumntype{Y}{>{\RaggedRight\arraybackslash}X}

\definecolor{milcMagenta}{HTML}{E6007E}
\definecolor{milcVino}{HTML}{5A0038}
\definecolor{milcTurquesa}{HTML}{0093A5}
\definecolor{milcVerde}{HTML}{58A923}
\definecolor{milcMostaza}{HTML}{E5B400}
\definecolor{milcOcre}{HTML}{B86600}
\definecolor{milcNegro}{HTML}{241816}
\definecolor{milcGris}{HTML}{F7F7F4}
\definecolor{milcLinea}{HTML}{E8E2DD}
\definecolor{milcGradePrimary}{HTML}{84CC16}
\definecolor{milcGradeDark}{HTML}{4D7C0F}
\definecolor{milcGradeSoft}{HTML}{ECFCCB}
\definecolor{milcGradeAccent}{HTML}{CFE7FF}
\definecolor{milcGradeText}{HTML}{FFFFFF}
\color{milcNegro}

\pagestyle{fancy}
\fancyhf{}
\lhead{\small Tecnología e Informática · Bebras}
\rhead{\small Momento 2 · MILC v3}
\cfoot{\small\thepage}
\renewcommand{\headrulewidth}{0pt}

\newtcolorbox{titlebox}[1]{
  enhanced,colback=#1,colframe=#1,arc=7mm,boxrule=0pt,
  left=7mm,right=7mm,top=3mm,bottom=3mm,
  before skip=8pt,after skip=8pt,
  fontupper=\sffamily\bfseries\Large\color{white}
}

\newtcolorbox{softbox}[3]{
  enhanced,colback=#2,colframe=#1,borderline west={4pt}{0pt}{#1},
  arc=2mm,boxrule=.4pt,left=4mm,right=4mm,top=3mm,bottom=3mm,
  before skip=6pt,after skip=8pt,title={#3},coltitle=white,
  colbacktitle=#1,fonttitle=\sffamily\bfseries,fontupper=\RaggedRight,
  attach boxed title to top left={xshift=3mm,yshift=-2mm},
  boxed title style={arc=2mm, boxrule=0pt}
}

\newtcolorbox{drawbox}[2]{
  enhanced,colback=white,colframe=#1,arc=4mm,boxrule=1pt,
  left=5mm,right=5mm,top=4mm,bottom=4mm,before skip=8pt,after skip=8pt,
  title={#2},coltitle=white,colbacktitle=#1,fonttitle=\sffamily\bfseries,
  fontupper=\RaggedRight,
  attach boxed title to top left={xshift=4mm,yshift=-2mm},
  boxed title style={arc=3mm, boxrule=0pt}
}

\newtcolorbox{infoband}[1]{
  enhanced,colback=#1!8,colframe=#1!50,arc=2mm,boxrule=.5pt,
  left=4mm,right=4mm,top=2mm,bottom=2mm,before skip=4pt,after skip=4pt,
  fontupper=\small\RaggedRight
}

% Puente narrativo entre secciones (contrato editorial Conéctate).
% Da continuidad: "Acabas de X. Ahora vas a Y porque Z."
% Visualmente: banda izquierda en color del grado, texto en itálica pequeña.
\newtcolorbox{puentebox}{
  enhanced,colback=milcGradeSoft,colframe=milcGradeDark,
  borderline west={3pt}{0pt}{milcGradeDark},
  arc=0mm,boxrule=0pt,left=5mm,right=4mm,top=2.5mm,bottom=2.5mm,
  before skip=4pt,after skip=8pt,
  fontupper=\itshape\small\color{milcNegro}\RaggedRight
}
% La flecha va en modo matematico a proposito: Helvetica Neue NO tiene el glifo
% U+2192, asi que \textrightarrow se compilaba a NADA («Missing character: There
% is no → in font Helvetica Neue») y el puente salia sin flecha. En modo
% matematico la flecha viene de la fuente matematica y si se dibuja.
\newcommand{\puente}[1]{%
  \begin{puentebox}%
    \textcolor{milcGradeDark}{$\rightarrow$}\hspace{2mm}#1%
  \end{puentebox}%
}

\newcommand{\linead}[1]{\noindent\color{milcLinea}\rule{#1}{0.5pt}\color{milcNegro}}
\newcommand{\respuesta}[1]{\par\vspace{2mm}\foreach \n in {1,...,#1}{\noindent\color{milcLinea}\rule{\linewidth}{0.5pt}\par\vspace{5mm}}\color{milcNegro}}
\newcommand{\checkbox}{\textcolor{milcGradeDark}{\fbox{\phantom{x}}}\hspace{5pt}}

\newcommand{\phasecard}[4]{
\begin{tcolorbox}[enhanced,colback=#3,colframe=#2,arc=3mm,boxrule=.5pt,height=3.05cm,valign=center,halign=center,left=1.5mm,right=1.5mm,top=2mm,bottom=2mm]
{\sffamily\bfseries\fontsize{22}{24}\selectfont\textcolor{#2}{#1}}\par
{\sffamily\bfseries\small #4}
\end{tcolorbox}
}

% ── Piezas del rediseno 2026 (legibilidad 6.º-11.º) ───────────────────
% Banda de estacion: reemplaza el titlebox macizo. Titulo a la izquierda,
% dato practico (tiempo/modalidad) a la derecha, en una sola linea.
\newcommand{\milcestacion}[2]{%
\begin{tcolorbox}[enhanced,colback=#1,colframe=#1,arc=2mm,boxrule=0pt,
  left=5mm,right=5mm,top=2.6mm,bottom=2.6mm,before skip=11pt,after skip=7pt]
{\sffamily\bfseries\fontsize{15}{18}\selectfont\color{white}#2}
\end{tcolorbox}}

% Tarjeta de paso numerada: sustituye las tablas «Pilar / Aplicacion», que
% eran formato de consulta docente, no de estudiante. El numero va en un
% circulo sobre el borde para que la secuencia se lea de un vistazo.
\newcommand{\milcpaso}[3]{%
\begin{tcolorbox}[enhanced,colback=white,colframe=milcLinea,arc=1.5mm,boxrule=.9pt,
  left=14mm,right=4mm,top=3mm,bottom=3mm,before skip=4pt,after skip=4pt,
  fontupper=\RaggedRight,
  overlay={\node[anchor=north west,circle,fill=#2,text=white,inner sep=0pt,
    minimum size=7.4mm,font=\sffamily\bfseries\small]
    at ([xshift=3.6mm,yshift=-3.2mm]frame.north west) {#1};}]
#3
\end{tcolorbox}}

% Casilla marcable de tamano util con lapiz (la anterior era \fbox{\phantom{x}},
% de alto variable segun el texto que la seguia).
\newcommand{\milcmarca}{\framebox[3.8mm]{\rule{0pt}{3.4mm}}}

% Fila marcable de lista suelta (checks de la estacion 1).
\newcommand{\milccheckrow}[1]{%
\par\vspace{1.8mm}\noindent
\makebox[6.4mm][l]{\raisebox{-.55mm}{\textcolor{milcNegro}{\milcmarca}}}%
\parbox[t]{\dimexpr\linewidth-6.4mm\relax}{\RaggedRight #1}\par}

% Celda marcable dentro de la rubrica (tabla de autoevaluacion). Misma
% sangria francesa, pero pensada para vivir dentro de una columna X.
\newcommand{\milcopcion}[1]{%
\makebox[5.2mm][l]{\raisebox{-.55mm}{\textcolor{milcNegro}{\milcmarca}}}%
\parbox[t]{\dimexpr\linewidth-5.2mm\relax}{\RaggedRight #1}}

% ── Recursos visuales de la guía ──────────────────────────────────────
% Declarados en el bloque `recursos:` del YAML y emitidos por el builder.
% \guiaFigura   → imagen rasterizada o PDF (foto, carta celeste, montaje).
% \guiaDiagrama → fuente TikZ/circuitikz (.tex) incrustada como vector:
%                 es la vía para circuitos y esquemáticos editables.
% #1 = ruta relativa al .tex · #2 = pie (puede ir vacío)
\newcommand{\guiaFigura}[2]{%
\begin{center}
\includegraphics[width=0.88\linewidth,height=0.38\textheight,keepaspectratio]{#1}\\[1.5mm]
{\footnotesize\itshape\textcolor{milcNegro!75}{#2}}
\end{center}
}

\newcommand{\guiaDiagrama}[2]{%
\begin{center}
\input{#1}\\[1.5mm]
{\footnotesize\itshape\textcolor{milcNegro!75}{#2}}
\end{center}
}

\begin{document}
\thispagestyle{empty}

% ── Portada ───────────────────────────────────────────────────────────
% Antes: un rectangulo de color a pagina completa. Se veia bien en pantalla,
% pero en la fotocopiadora del colegio se comia el toner y salia gris sucio.
% Ahora la banda ocupa el tercio superior y el resto va en blanco.
\begin{tikzpicture}[remember picture,overlay]
  \path[fill=milcGradePrimary]
    (current page.north west) rectangle ([yshift=-10.6cm]current page.north east);

  \node[anchor=north west,text=milcGradeText] at ([xshift=2.0cm,yshift=-1.55cm]current page.north west)
    {\sffamily\bfseries\fontsize{12}{14}\selectfont MOMENTO};
  \node[anchor=north west,text=milcGradeText,opacity=.85] at ([xshift=2.0cm,yshift=-2.35cm]current page.north west)
    {\sffamily\bfseries\fontsize{8.5}{10}\selectfont BEBRAS · PENSAR COMO UNA MÁQUINA};
  \node[anchor=north east,text=milcGradeText,opacity=.85,align=right] at ([xshift=-2.0cm,yshift=-1.55cm]current page.north east)
    {\sffamily\bfseries\fontsize{9}{12}\selectfont I.E. Sor María Juliana\\Cartago, Valle del Cauca};

  \node[anchor=west,text=milcGradeText] (gradonum) at ([xshift=1.85cm,yshift=-7.1cm]current page.north west)
    {\sffamily\bfseries\fontsize{112}{112}\selectfont 2};
  % El circulito de grado (°) solo aplica a las guias de grado («11°»); en
  % Bebras y Semillero el numero grande es un momento/modulo, no un grado.
  % Se posiciona relativo al nodo `gradonum`, no en coordenadas absolutas.
  

  \node[anchor=west,text=milcGradeText,text width=11.4cm,align=left] at ([xshift=1.5cm,yshift=0.5cm]gradonum.east)
    {
     {\sffamily\bfseries\fontsize{9}{11}\selectfont\MakeUppercase{Curso «Pensar como una máquina» · Pensamiento computacional}}\\[3mm]
     {\sffamily\bfseries\fontsize{21}{25}\selectfont\setlength{\baselineskip}{27pt}Cazador de patrones --- ver lo que se repite\\[4pt]para predecir lo que sigue}\\[3.5mm]
     {\sffamily\bfseries\fontsize{15}{18}\selectfont Bebras · Momento 2}
    };
\end{tikzpicture}

\vspace*{9.4cm}
\vfill

{\sffamily\bfseries\fontsize{8}{10}\selectfont\textcolor{milcGradeDark}{LO QUE TE LLEVAS HOY}}

\begin{tcolorbox}[enhanced,colback=white,colframe=milcNegro,arc=2mm,boxrule=1.6pt,
  left=5mm,right=5mm,top=4mm,bottom=4mm,before skip=3pt,after skip=12pt,
  fontupper=\RaggedRight\fontsize{11}{15.5}\selectfont]
Tu tabla «Tres patrones de mi entorno» con regla y predicción, más el simulacro A·B·C resuelto y explicado con el truco del residuo.
\end{tcolorbox}

\begin{tcolorbox}[enhanced,colback=milcGris,colframe=milcGris,arc=2mm,boxrule=0pt,
  left=5mm,right=5mm,top=3.5mm,bottom=3.5mm,after skip=12pt]
{\sffamily\bfseries\fontsize{8}{10}\selectfont\textcolor{milcNegro!55}{ESTA GUÍA ES DE}}\\[3mm]
\begin{tabularx}{\linewidth}{X p{3.2cm} p{3.2cm}}
\linead{\linewidth} & \linead{3.0cm} & \linead{3.0cm}\\[-1mm]
{\fontsize{8.5}{10}\selectfont\textcolor{milcNegro!55}{Nombre completo}} &
{\fontsize{8.5}{10}\selectfont\textcolor{milcNegro!55}{Grupo}} &
{\fontsize{8.5}{10}\selectfont\textcolor{milcNegro!55}{Fecha}}\\
\end{tabularx}
\end{tcolorbox}

\vfill

\noindent\textcolor{milcLinea}{\rule{\linewidth}{1pt}}\\[2mm]
{\sffamily\bfseries\fontsize{9.5}{12}\selectfont Autor: Dr. Álvaro Cárdenas Orozco}\\[1mm]
{\sffamily\fontsize{8.5}{11}\selectfont\textcolor{milcNegro!55}{Metodología MILC · apertura ancestral · pensamiento computacional · triángulo Dussel–estoicismo–Floridi}}

\newpage

\section*{Apertura: Cazador de patrones: ver lo que se repite para predecir lo que sigue}

\begin{softbox}{milcGradeDark}{milcGradeSoft}{Saber ancestral}
En el Pacífico colombiano, cuando suena el \textbf{currulao}, la marimba de chonta no toca cualquier cosa: teje una célula rítmica que vuelve una y otra vez. El cununo y el bombo marcan un golpe que se repite ---\emph{pum, pa, pum-pum, pa}--- como un tejido sonoro que no se rompe. Sobre esa base predecible, el cantador improvisa y los bailarines responden. ¿Y cómo saben los bailarines cuándo viene el golpe fuerte, cuándo zapatear, cuándo girar el pañuelo? No lo adivinan: \textbf{reconocen el patrón}. Su oído ya aprendió lo que se repite, así que pueden anticipar lo que sigue. Lo mismo pasa con los tejidos del Pacífico y de La Guajira: una mochila wayuu o una canasta en werregue repiten un diseño ---rombo, línea, rombo, línea--- que la tejedora lleva en la memoria. Reconocer lo que se repite, en el sonido o en la fibra, \textbf{es ya un acto de pensamiento matemático}.
\end{softbox}

\begin{softbox}{milcTurquesa}{milcGris}{Saber contemporáneo}
\textbf{Reconocer patrones} es la segunda pieza del pensamiento computacional: encontrar regularidades, repeticiones y secuencias en un conjunto de datos. Un patrón es cualquier cosa que vuelve a aparecer siguiendo una \textbf{regla}: la tabla del 5 siempre termina en 0 o en 5 porque la regla es ``suma 5 cada vez''. Una vez que ves la regla, ganas dos superpoderes: \textbf{predecir} lo que sigue sin verlo (como saber a qué hora pasa el bus de la ruta 2 sin mirar el horario) y \textbf{generalizar} (resolver la regla una vez y aplicarla a mil casos, no solo a diez). En las tareas Bebras, el reconocimiento de patrones aparece en secuencias numéricas, ciclos de colores y diseños repetidos. La clave no es dibujar cada paso: es hallar la regla y \emph{usarla}.
\end{softbox}

\begin{softbox}{milcMostaza}{milcGris}{Pregunta-puente}
\textit{Quien baila currulao ``sabe'' cuándo viene el golpe porque reconoce el patrón rítmico. ¿Y si te dijera que esa misma habilidad ---ver lo que se repite para predecir lo que sigue--- es justo lo que te pide una tarea Bebras?}
\end{softbox}

\begin{softbox}{milcVerde}{milcGris}{Saber y saber hacer}
Al terminar podrás: (1) \textbf{identificar} lo que se repite en una secuencia de números, colores o figuras, nombrando la regla que la genera; (2) \textbf{explicar} con tus palabras la diferencia entre repetición, regla, predicción y generalización, con ejemplos de tu entorno; (3) \textbf{aplicar} el truco del residuo para hallar una posición lejana en un patrón cíclico, sin dibujar todos los casos; (4) resolver un simulacro A·B·C de patrones justificando cada respuesta.
\end{softbox}



\small
\begin{tabularx}{\textwidth}{>{\bfseries}p{3.0cm}Y}
\rowcolor{milcGradePrimary!12}Autor & Dr. Álvaro Cárdenas Orozco\\
DBA & Desarrolla el pensamiento computacional --- descomposición, patrones, abstracción y algoritmos --- para resolver problemas (transversal 6°--11°, alineado a los lineamientos MEN de Tecnología e Informática)\\
Referentes & Valentina Dagienė (Bebras) · Pensamiento computacional (Wing) · Dussel · Estoicismo · Floridi · Saberes del Valle y el Pacífico\\
Producto final & Tu tabla «Tres patrones de mi entorno» con regla y predicción, más el simulacro A·B·C resuelto y explicado con el truco del residuo.\\
\end{tabularx}
\normalsize

\puente{Ya sabes que el currulao predice el golpe porque reconoce el patrón. Ahora mira el mapa de la sesión: vas a pasar de ver lo que se repite a usarlo para predecir, paso por paso.}

\milcestacion{milcGradeDark}{Ruta MILC: así vamos a aprender}
\begin{tabularx}{\textwidth}{>{\centering\arraybackslash}X>{\centering\arraybackslash}X>{\centering\arraybackslash}X>{\centering\arraybackslash}X}
\phasecard{1}{milcVerde}{milcGris}{Escucha\\[-1mm]\small Observo, escucho y pregunto.} &
\phasecard{2}{milcTurquesa}{milcGris}{Entender\\[-1mm]\small Sistematizo y ordeno.} &
\phasecard{3}{milcMagenta}{milcGris}{Hacer\\[-1mm]\small Construyo y aplico.} &
\phasecard{4}{milcVino}{milcGris}{Cierre\\[-1mm]\small Evalúo y reflexiono.}
\end{tabularx}


\puente{Antes de nombrar qué es un patrón, conviene buscarlo en lo que ya conoces. Recorre tu casa, tu barrio y tu colegio: los patrones están ahí, esperando que los veas.}

\milcestacion{milcVerde}{Estación 1 de 3 · Escucha}

\begin{softbox}{milcVerde}{milcGris}{Escena para escuchar}
\textbf{Recorre con la mirada tu casa, tu barrio y tu colegio} buscando cosas que se repitan: baldosas que alternan color, rejas con el mismo diseño, el horario del bus de la ruta 2, los días de mercado en Cartago, el diseño de una mochila o una ruana. Elige \textbf{tres patrones distintos}. Para cada uno, anota qué elemento se repite y cuál es la regla que lo genera ---en tus palabras, sin tecnicismos.
\end{softbox}

{\sffamily\bfseries\fontsize{8}{10}\selectfont\textcolor{milcNegro!55}{MARCA CUANDO LO HAYAS HECHO}}
\milccheckrow{Encontraste 3 patrones reales de tu entorno (no del texto): uno puede ser visual, uno numérico o de tiempo, uno de otro tipo.}
\milccheckrow{Para cada patrón describiste qué elemento se repite con claridad suficiente para que un compañero lo entienda.}
\milccheckrow{Para al menos uno de los tres escribiste cómo usarías la regla para predecir algo útil (por ejemplo, a qué hora pasa el bus).}

\begin{infoband}{milcVerde}


\textbf{En tu cuaderno (Actividad 1 · IDENTIFICA).} Título: «Actividad 1 --- Tres patrones de mi entorno». \textbf{Formato:} tabla de 3 columnas: \emph{dónde lo vi} · \emph{qué se repite} · \emph{regla en mis palabras}. \textbf{Extensión:} 3 filas completas + una frase final que diga cuál de los tres patrones te sirve para predecir algo útil. \textbf{Sabes que terminaste cuando} tu tabla tiene los tres patrones, cada uno con su regla escrita, y la frase final está redactada.
\end{infoband}

\puente{Acabas de cazar patrones en tu entorno. Ahora vamos a ponerles nombre y a entender las cuatro ideas que hacen que un patrón sea útil: repetición, regla, predicción y generalización.}

\milcestacion{milcTurquesa}{Estación 2 de 3 · Entender}

\begin{softbox}{milcTurquesa}{milcGris}{Lo que vas a entender}
Reconocer patrones no es memorizar: es descubrir la regla que genera la secuencia. Esa regla te da dos poderes: \textbf{predecir} (saber lo que sigue sin verlo) y \textbf{generalizar} (resolver una vez y aplicar a mil casos). Aquí están las cuatro ideas que hacen que un patrón sea útil, con ejemplos de Cartago y el Pacífico.
\end{softbox}

\milcpaso{1}{milcTurquesa}{\textbf{Repetición:} algo vuelve a aparecer igual, una y otra vez. El cununo del currulao repite su golpe; la mochila wayuu repite el rombo. Si lo notas, ya no tienes que mirar cada caso por separado: el patrón hace el trabajo.}
\milcpaso{2}{milcTurquesa}{\textbf{La regla:} es la instrucción oculta que genera la secuencia. En la tabla del 5 (5, 10, 15, 20\ldots) la regla es ``suma 5 cada vez''; por eso siempre termina en 0 o en 5. Hallar la regla es el corazón del trabajo: con ella predices cualquier término.}
\milcpaso{3}{milcTurquesa}{\textbf{Predecir:} si conoces la regla, sabes lo que sigue sin esperar a verlo. El bus de la ruta 2 pasa a los minutos 00 y 30: ya sabes a qué hora salir de casa sin mirar ningún horario. En Bebras, predecir te ahorra dibujarlo todo.}
\milcpaso{4}{milcTurquesa}{\textbf{Generalizar:} en vez de resolver mil casos uno por uno, resuelves la regla una sola vez y la aplicas a todos. El \textbf{truco del residuo}: si el ciclo mide 5 y quieres la posición 23, divide 23 entre 5 --- caben 4 grupos (20) y sobran 3 --- la posición 3 del ciclo es tu respuesta.}

\begin{drawbox}{milcTurquesa}{Cómo hallar la regla de un patrón}
\textbf{1. Mira los primeros términos} y pregúntate: ¿qué cambia de uno al otro? \\[1mm] \textbf{2. Prueba una operación simple:} ¿suma, resta, multiplicación, ciclo? \\[1mm] \textbf{3. Comprueba} que tu regla funciona en todos los términos que ya ves. \\[1mm] \textbf{4. Aplica la regla} para predecir el siguiente término sin dibujarlo. \\[1mm] \textbf{5. Para posiciones lejanas,} usa el truco del residuo: divide la posición entre el tamaño del ciclo y usa el resto. Si el resto es 0, es el último del ciclo.
\end{drawbox}

\begin{softbox}{milcMostaza}{milcGris}{Errores comunes y su corrección}
\textbf{Copiar la secuencia sin buscar la regla:} el problema no pide que dibujes los 23 pasos, sino que halles la regla y la apliques. \textbf{Confundir el ciclo con la posición:} el ciclo tiene longitud fija (por ejemplo, 5), pero la posición puede ser cualquier número. \textbf{Olvidar que el residuo 0 es el último del ciclo:} si divides 20 entre 5 y el resto es 0, la cuenta es la 5.ª del patrón, no la 0.ª.
\end{softbox}

\begin{infoband}{milcTurquesa}


\textbf{En tu cuaderno (Actividad 2 · EXPLICA).} Título: «Actividad 2 --- La regla de mi patrón». \textbf{Formato:} párrafo corto en lenguaje sencillo. \textbf{Extensión:} 5 a 7 renglones. \textbf{Sabes que terminaste cuando} explicaste la regla de uno de tus tres patrones con tus palabras, mostraste cómo predecir el siguiente elemento y conectaste tu patrón con el ritmo del currulao o un saber de tu familia.
\end{infoband}

\puente{Ya distingues la regla de la secuencia. Es hora de usarla en retos reales al estilo Bebras, donde lo que importa es hallar la regla antes de responder.}

\milcestacion{milcMagenta}{Estación 3 de 3 · Hacer}

\begin{softbox}{milcMagenta}{milcGris}{Cómo vas a construir tu producto}
La regla no es adorno: se usa. Vas a resolver tres retos al estilo Bebras sobre reconocimiento de patrones, de fácil (A) a difícil (C). Lee con calma --- la clave está en hallar la regla antes de responder, no en adivinar.
\end{softbox}

\milcpaso{1}{milcMagenta}{\textbf{Descomposición:} parte el enunciado en datos y pregunta. ¿Qué patrón te dan? ¿Qué posición te piden?}
\milcpaso{2}{milcMagenta}{\textbf{Patrones:} identifica el ciclo y su longitud. ¿Cuántos elementos tiene antes de repetirse?}
\milcpaso{3}{milcMagenta}{\textbf{Abstracción:} ignora los datos de adorno; quédate solo con el ciclo y la posición pedida.}
\milcpaso{4}{milcMagenta}{\textbf{Algoritmo:} aplica el truco del residuo --- divide la posición entre la longitud del ciclo y usa el resto para ubicarte dentro del patrón.}

\begin{drawbox}{milcMagenta}{El reto: la tejedora wayuu y la cuenta 23}
\checkbox \textbf{Reto (Nivel C).} Una tejedora wayuu arma un collar repitiendo \emph{siempre} este patrón de 5 cuentas en este orden: 1.~roja · 2.~negra · 3.~verde · 4.~amarilla · 5.~azul, y vuelve a empezar. ¿De qué color es la cuenta número~23? \\[2mm]
\checkbox Marca: \quad \fbox{\phantom{x}}~Roja \quad \fbox{\phantom{x}}~Negra \quad \fbox{\phantom{x}}~Verde \quad \fbox{\phantom{x}}~Amarilla \quad \fbox{\phantom{x}}~Azul \\[2mm]
\checkbox Escribe el \textbf{truco que usaste} (residuo) y explica por qué funciona sin dibujar las 23 cuentas.
\end{drawbox}

\begin{softbox}{milcMagenta}{milcGris}{Plantilla de trabajo}
\textbf{Resuélvelo paso a paso (truco del residuo).} El patrón tiene \textbf{5 cuentas} que se repiten. Quieres la posición 23. Divide: $23 \div 5 = 4$ grupos completos con \textbf{residuo 3}. El residuo 3 corresponde a la \textbf{3.ª cuenta del patrón}, que es \textbf{verde}. La respuesta es \textbf{verde}. ¿Por qué funciona? Porque 4 grupos completos (20 cuentas) ya se «cerraron»; las 3 que sobran son el inicio del quinto grupo: roja, negra, \textbf{verde}. No necesitabas dibujar las 23 cuentas.
\end{softbox}

\begin{infoband}{milcMagenta}


\textbf{En tu cuaderno (Actividad 3 · APLICA).} Título: «Actividad 3 --- El collar de la tejedora wayuu». \textbf{Formato:} respuesta marcada + 4 a 5 renglones de explicación con el truco del residuo. \textbf{Extensión:} la operación de división, el residuo, la posición en el patrón y el color. \textbf{Sabes que terminaste cuando} tu explicación permite a alguien verificar la respuesta \emph{sin dibujar las 23 cuentas}.
\end{infoband}

\puente{Resolviste el simulacro. Ahora deja tu evidencia: la tabla de patrones y las soluciones explicadas, para mostrar que entiendes por qué funciona cada respuesta.}

\milcestacion{milcMostaza}{Producto de la sesión}
\textbf{Tabla «Tres patrones de mi entorno» + simulacro A·B·C resuelto con el truco del residuo explicado}.
\begin{softbox}{milcMostaza}{milcGris}{Criterios de calidad}
(1) Tu tabla tiene 3 patrones reales y distintos, cada uno con la regla escrita en tus palabras y una predicción útil. (2) En el reto A (manilla) identificaste el ciclo de 2 colores y marcaste la respuesta correcta (verde). (3) En el reto B (tambor del currulao) aplicaste la regla ``multiplica por 2'' y obtuviste 48. (4) En el reto C (collar wayuu) usaste el truco del residuo ($23 \div 5 = 4$ r.~3 $\to$ verde) y lo explicaste sin dibujar las 23 cuentas. (5) Tu trabajo muestra que entendiste que generalizar es resolver la regla una vez y aplicarla a cualquier posición, no contar caso por caso.
\end{softbox}

\puente{Terminaste el producto. Detente a mirar qué cambió en ti --- no la nota, sino tu manera de ver el mundo: ¿ya ves reglas donde antes solo veías cosas sueltas?}

\milcestacion{milcVino}{Evaluación liberadora · cinco dimensiones}

\begin{softbox}{milcVino}{milcGris}{Sentido de la evaluación}
Evaluar no es solo revisar una nota. Es reconocer qué aprendí, qué cambió en mí, qué emociones manejé, cómo me relaciono con la ciudad y la comunidad, y qué le dejaré a quien viene detrás.
\end{softbox}

\textbf{Mi reflexión liberadora en cinco dimensiones.} Completa cada frase iniciada:

\small
\begin{tabularx}{\textwidth}{>{\bfseries\RaggedRight\arraybackslash}p{3.4cm}Y>{\RaggedRight\arraybackslash}p{4.6cm}}
\rowcolor{milcVino}\color{white}Dimensión & \color{white}\bfseries Mi respuesta (frame starter) & \color{white}\bfseries Algo que descubrí\\
1. Personal & Lo que más cambió en mí fue \linead{6.5cm} & \linead{4.2cm}\\
2. Emocional & La emoción más fuerte fue \linead{2.5cm} y la regulé \linead{3cm} & \linead{4.2cm}\\
3. Ciudadana & En democracia esto importa porque \linead{6cm} & \linead{4.2cm}\\
4. Local (Cartago) & En mi barrio sería útil para \linead{6.5cm} & \linead{4.2cm}\\
5. Intergeneracional & A mi abuela/abuelo le diría \linead{6.5cm} & \linead{4.2cm}\\
\end{tabularx}
\normalsize

\begin{infoband}{milcVino}
\textbf{Para la próxima sustentación.} Vuelve a esta tabla en seis meses. Verás que las respuestas cambian — no porque lo aprendido cambie, sino porque tú cambias. La evaluación liberadora es una conversación contigo a través del tiempo.
\end{infoband}


\puente{Cerramos pensando en grande: tres voces para entender qué significa, para ti y para tu comunidad, aprender a ver lo que se repite.}

\milcestacion{milcGradeDark}{Triángulo de pensamiento · cierre filosófico}

\begin{softbox}{milcVino}{milcGris}{Dussel · lente del nosotros}
\textit{Los saberes del pueblo también son ciencia.}

Los patrones del currulao, de la mochila wayuu y de la canasta en werregue son conocimiento matemático del Sur: ritmo y geometría que nuestros pueblos pensaron mucho antes que cualquier manual escolar. El reconocimiento de patrones no llegó de afuera --- tu comunidad ya lo practicaba en el tejido y en la música.

\textbf{¿Qué patrones de mi cultura ---en la música, el tejido o la cocina--- ya eran matemática, aunque nadie los llamara así?}
\end{softbox}
\linead{\linewidth}\\[1mm]
\linead{\linewidth}

\begin{softbox}{milcMostaza}{milcGris}{Estoicismo · Marco Aurelio · cuidado interior}
\textit{Observa el curso regular de las cosas: todo lo que sucede vuelve a suceder, y quien lo comprende deja de turbarse por ello.}

Cuando algo parece impredecible y el tiempo corre --- un examen, una situación difícil ---, busca el patrón: ¿qué se repite aquí? Reconocer la regularidad no elimina el problema, pero sí te devuelve el control sobre cómo lo enfrentas.

\textbf{¿Qué patrones reconozco en mis propios hábitos ---cuándo me distraigo, cuándo rindo--- y qué podría predecir y mejorar en mí?}
\end{softbox}
\linead{\linewidth}\\[1mm]
\linead{\linewidth}

\begin{softbox}{milcTurquesa}{milcGris}{Floridi · lente de la infoesfera}
\textit{Los algoritmos que nos rodean viven de patrones: aprenden lo que se repite en nuestros datos para anticipar lo que haremos después.}

Las aplicaciones que usas cada día reconocen patrones en ti: qué ves, a qué hora, qué compras, con quién hablas. Aprender a ver patrones es también aprender a entender cómo te leen las máquinas --- y a decidir qué muestras.

\textbf{¿Qué patrones míos crees que ya conocen las apps que uso, y para qué los usan?}
\end{softbox}
\linead{\linewidth}\\[1mm]
\linead{\linewidth}

\begin{infoband}{milcNegro}
\textbf{Nota para el docente.} Las tres formulaciones del triángulo son la síntesis del autor de la guía, no citas textuales de los pensadores. Si vas a citarlos en clase, remite a la obra original.
\end{infoband}

\milcestacion{milcVino}{Mi autoevaluación · marca una casilla por fila}

{\small
\setlength{\extrarowheight}{2.2pt}
\begin{tabularx}{\textwidth}{>{\RaggedRight\arraybackslash\bfseries}p{3.0cm}YYY}
\rowcolor{milcVino}\color{white}\bfseries Criterio & \color{white}\bfseries Lo logré & \color{white}\bfseries En proceso & \color{white}\bfseries Necesito apoyo\\[.6mm]
Comprensión del tema & \milcopcion{Explico las ideas con mis palabras.} & \milcopcion{Reconozco las ideas, falta precisión.} & \milcopcion{Necesito repasar lo básico.}\\[.8mm]
\rowcolor{milcGris}Pensamiento computacional & \milcopcion{Usé los pilares al armar mi producto.} & \milcopcion{Usé algunos pilares.} & \milcopcion{Necesito releer la estación 2.}\\[.8mm]
Aplicación y producto & \milcopcion{Mi producto cumple los criterios.} & \milcopcion{Está armado, con ajustes pendientes.} & \milcopcion{Aún está en construcción.}\\[.8mm]
\rowcolor{milcGris}Reflexión 5 dimensiones & \milcopcion{Respondí cada una con honestidad.} & \milcopcion{Respondí algunas con detalle.} & \milcopcion{Mis respuestas son muy generales.}\\[.8mm]
Triángulo de pensamiento & \milcopcion{Conecté las 3 voces con mi trabajo.} & \milcopcion{Conecté al menos una.} & \milcopcion{No las relaciono con mi proceso.}\\[.8mm]
\end{tabularx}}

\vspace{3mm}
\textbf{Hoy entendí que:}\\[1mm]
\linead{\linewidth}\\[1mm]
\linead{\linewidth}

\textbf{Mi compromiso para la próxima sesión es:}\\[1mm]
\linead{\linewidth}\\[1mm]
\linead{\linewidth}

\begin{softbox}{milcMostaza}{milcGris}{Cita de despedida · Marco Aurelio}
\textit{``El obstáculo en el camino se vuelve el camino. Lo que estorba al camino se vuelve camino.''}\\[1mm]
Cada error de hoy, bien observado, se convierte en pieza del camino que pisarás mañana.
\end{softbox}

\small
\begin{tabularx}{\textwidth}{>{\bfseries}p{3.6cm}Y}
\rowcolor{milcGradeDark!12}Nombre completo: & \linead{10cm}\\
Fecha: & \linead{4cm} \quad Curso/grupo: \linead{4cm}\\
Mi firma: & \linead{9cm}\\
\end{tabularx}
\normalsize

\begin{infoband}{milcGradeDark}
\textbf{Para la próxima sesión.} Elige UNA secuencia numérica que hayas visto esta semana --- en el precio de algo, en el horario del bus, en los números de las casas de tu cuadra --- y escribe: cuál es la regla, cómo predecirías el siguiente elemento y si el truco del residuo sirve para encontrar un término lejano. Trae el ejemplo para compartir con el grupo.
\end{infoband}

\end{document}
